\section*{Chapter \ref{ch:mx:fragmentation-and-routing} --- Fragmentation and Routing}

\begin{solution}{exo:mx:fragmentation-and-routing:1}
Bid 10.01 for 500 (V2; V3's 50 shares are an odd lot, not protected); offer 10.02 for 1\,600 (V1 200, V3 400, V4 1\,000).
\end{solution}

\begin{solution}{exo:mx:fragmentation-and-routing:2}
\omterm{def:mx:fragmentation-and-routing:locked}{Locked}, then \omterm{def:mx:fragmentation-and-routing:locked}{crossed}. In the lock there is nothing to earn at equal prices, before fees; in the cross an arbitrageur buys on B at 10.02 and sells on A at
10.03, one cent a share, if both quotes are still there when its orders arrive.
\end{solution}

\begin{solution}{exo:mx:fragmentation-and-routing:3}
Each change is seen one millisecond late, so at most $500\times0.001=50\%$ of the time. The simulated bound is $0.44\times0.0005=0.022\%$; the measured
0.0098\% is lower because changes come in bursts whose late windows overlap.
\end{solution}

\begin{solution}{exo:mx:fragmentation-and-routing:4}
$\psi\Omega\psi'=0.76$. Ordering (1, 2): $\psi F=(0.8,\,0.346)$, shares 0.842 and 0.158. Ordering (2, 1): shares 0.645 for market 2 and 0.355 for market 1.
Bounds: market 1 between 0.355 and 0.842, market 2 between 0.158 and 0.645.
\end{solution}

\begin{solution}{exo:mx:fragmentation-and-routing:5}
$\alpha_\perp\propto(0.3,0.2)$, since $\alpha'\alpha_\perp=-0.06+0.06=0$. \omterm{def:mx:fragmentation-and-routing:cs}{Component shares} 0.6 and 0.4: market 1 adjusts less to the gap and leads.
\end{solution}

\begin{solution}{exo:mx:fragmentation-and-routing:6}
\omterm{def:mx:fragmentation-and-routing:cs}{Component shares} are normalised weights of $\alpha_\perp$ and nothing forces the entries to share a sign. A negative weight comes from a market whose
adjustment coefficient has the sign of a market that overshoots: C's quotes, ten seconds late, move after the others and by more than the gap, so the
common component loads against them. The \omterm{def:mx:fragmentation-and-routing:is}{information share}, a variance share, cannot be negative.
\end{solution}

\begin{solution}{exo:mx:fragmentation-and-routing:7}
One-second mids and five lags give B 0.37 to 0.59, A 0.23 to 0.44, C 0.13 to 0.28: B comes first. The delays of 0.2 and 1 second are shorter than the
time each venue's quotes take to follow $P^\ast$ (seconds to tens of seconds in \texttt{firm.tape}), so the head start is lost in each venue's own quote
noise; with delays of 2 and 10 seconds A is found. Finer sampling helps only if the quotes react faster than the sampling interval.
\end{solution}

\begin{solution}{exo:mx:fragmentation-and-routing:8}
Volume is not \omterm{def:mx:fragmentation-and-routing:discovery}{price discovery}. In the simulation C's volume rose from 128\,100 to 230\,900 shares when an arbitrageur began trading against its stale
quotes, while its \omterm{def:mx:fragmentation-and-routing:is}{information share} stayed between 0.01 and 0.04: the extra volume was other traders taking C's late prices. Estimate \omterm{def:mx:fragmentation-and-routing:is}{information shares}
before and after, not volume shares.
\end{solution}

\begin{solution}{pb:mx:fragmentation-and-routing:1}
\textbf{1.} One security trading on several venues at once; competition on fees, speed and features, against split liquidity and routing.
\textbf{2.} Three \texttt{firm.tape} flows on one efficient price, seen at once on A, 2 seconds late on B, 10 on C, with 70\% and 45\% of A's activity on B
and C.
\textbf{3.} 38\%, 30\% and 32\% on A, B and C, time-weighted.
\textbf{4.} Lower transaction costs and faster executions, more short-term volatility, prices closer to a random walk.
\textbf{5.} Each venue's top of book reaches it after 0.5 ms; it computes the NBBO from round-lot quotes and publishes 20 microseconds later.
\textbf{6.} 0.0098\%; 0.002\% at 50 microseconds, 0.018\% at 1 ms, 0.084\% at 5, 0.33\% at 20, 1.5\% at 100.
\textbf{7.} 18.5\% of trades, nearly all the arbitrageur's; 0.02\% of the background's.
\textbf{8.} Dislocations several times a second in very active stocks, lasting one to two milliseconds.
\textbf{9.} Best bid on one venue equal to (\omterm{def:mx:fragmentation-and-routing:locked}{locked}) or above (\omterm{def:mx:fragmentation-and-routing:locked}{crossed}) the best offer on another.
\textbf{10.} \omterm{def:mx:fragmentation-and-routing:locked}{Locked} 34.4\% and \omterm{def:mx:fragmentation-and-routing:locked}{crossed} 10.0\% of the hour; 69 \omterm{def:mx:fragmentation-and-routing:locked}{crossed} episodes with a median of 4.8 seconds.
\textbf{11.} It removes crosses (565 episodes, median 41 microseconds, 0.07\% of the time) and leaves locks (42.6\%, median 0.44 seconds).
\textbf{12.} Mostly C's stale quotes (C's volume from 128\,100 to 230\,900): latency arbitrage.
\textbf{13.} To rescind Rule 611 and Rule 610(e), with a 60-day comment period.
\textbf{14.} $\Delta p_t=\alpha\beta'p_{t-1}+\sum\Gamma_k\Delta p_{t-k}+e_t$; $\mathrm{IS}_j=([\psi F]_j)^2/(\psi\Omega\psi')$.
\textbf{15.} A 0.63--0.79, B 0.19--0.34, C 0.01--0.04; \omterm{def:mx:fragmentation-and-routing:cs}{component shares} 0.75, 0.44 and $-0.20$.
\textbf{16.} The venues' innovations are correlated; the Cholesky factor attributes the common part to whichever market comes first.
\textbf{17.} B first (0.37--0.59), A second (0.23--0.44): the delays are shorter than the quotes' own reaction time and are lost in noise.
\textbf{18.} \emph{Named result}: with delays of 0, 2 and 10 seconds the \omterm{def:mx:fragmentation-and-routing:is}{information shares} are 0.63--0.79, 0.19--0.34 and 0.01--0.04 (\omterm{def:mx:fragmentation-and-routing:cs}{component shares}
0.75, 0.44, $-0.20$); the direct and consolidated best prices differ 0.002\% of the time at a 50-microsecond delay, 0.0098\% at half a millisecond,
0.084\% at 5 ms and 1.5\% at 100 ms, and at 18.5\% of trades, those of the fast arbitrageur.
\textbf{19.} No: C's volume nearly doubled with the arbitrageur while its \omterm{def:mx:fragmentation-and-routing:is}{information share} stayed near zero.
\textbf{20.} On the venue whose traders learn first, and in the arbitrageur's orders that carry it to the others.
\end{solution}

\begin{solution}{iq:mx:fragmentation-and-routing:1}
The best bid and best offer across protected quotations of all venues, computed by the securities information processor from the venues' data; a trader
with direct feeds computes it earlier and with its own timing, so the two differ whenever a change has not yet reached the processor.
\iqlookfor{\omterm{def:mx:fragmentation-and-routing:protected}{Protected quotes}; processor; latency.}
\end{solution}

\begin{solution}{iq:mx:fragmentation-and-routing:2}
A bid on one venue equal to the offer on another. US rules forbid displaying quotes that lock \omterm{def:mx:fragmentation-and-routing:protected}{protected quotes} elsewhere; a venue reprices or rejects the
order (or routes it to take the quote) because a lock sends contradictory signals and invites routing loops.
\iqlookfor{Definition; the rule; repricing.}
\end{solution}

\begin{solution}{iq:mx:fragmentation-and-routing:3}
Hasbrouck \omterm{def:mx:fragmentation-and-routing:is}{information shares} and Gonzalo--Granger \omterm{def:mx:fragmentation-and-routing:cs}{component shares} from a vector error-correction model on synchronised quotes. Pitfalls: wide bounds
when innovations are correlated (sample finer), a lead shorter than the sampling interval, stale quotes, and mistaking volume for discovery.
\iqlookfor{VECM; bounds; sampling frequency.}
\end{solution}

\begin{solution}{iq:mx:fragmentation-and-routing:4}
The venue, the venue's timestamp and my receive timestamp, the sequence number, and whether the quote is protected (round lot, displayed); I need them
to know how old each price is, to detect gaps, and to route to what is really there.
\iqlookfor{Timestamps; sequence; staleness.}
\end{solution}

\begin{solution}{iq:mx:fragmentation-and-routing:5}
Both effects exist: competition lowers fees and spreads, splitting liquidity raises search and routing costs. Look at matched stocks or natural
experiments (new venue entry, rule changes) and compare spreads, depth across venues, volatility and price efficiency.
\iqlookfor{Mechanisms on both sides; identification.}
\end{solution}

\begin{solution}{iq:mx:fragmentation-and-routing:6}
Anyone whose order or quote is checked or priced against the stale view: resting quotes on slow venues picked off by fast traders, marketable orders
routed to prices that have gone. Measure fills and mark-outs of trades that occur while the two views differ, as Bartlett and McCrary did with
microsecond timestamps, and the share of trades with a stale consolidated NBBO.
\iqlookfor{Who bears it; a measurement at trade times.}
\end{solution}
